\section{Certified finite-band boundary intervals}
\label{sec:intervals}

Fix a current face $S$ and a boundary vertex $v$.  Let
\begin{equation}
 a_v:=Q_{Sv},
 \qquad
 \chi_v^2:=a_v^\top Q_{SS}^{-1}a_v.
 \label{eq:boundary-leverage}
\end{equation}
The scalar $\chi_v$ is a response leverage.  It is closely related to the
energy needed to transmit a unit boundary demand through the active face.

\begin{lemma}[Energy-to-boundary interval]
\label{lem:boundary-interval}
Suppose $x_S$ is the exact restricted solution and $\widetilde x_S$ satisfies
\begin{equation}
 \|\widetilde x_S-x_S\|_{Q_{SS}}\leq\delta.
 \label{eq:energy-error}
\end{equation}
For every queried boundary vertex $v$,
\begin{equation}
 \left|
 Q_{vS}(\widetilde x_S-x_S)
 \right|
 \leq \chi_v\delta.
 \label{eq:boundary-error}
\end{equation}
Consequently a point estimate of any boundary demand affine in $Q_{vS}x_S$
becomes a certified interval after adding and subtracting $\chi_v\delta$.
\end{lemma}

\begin{proof}
Write the left side as
\[
 \left|
 (Q_{SS}^{-1/2}a_v)^\top
 Q_{SS}^{1/2}(\widetilde x_S-x_S)
 \right|.
\]
Cauchy--Schwarz and \cref{eq:boundary-leverage,eq:energy-error} prove the
claim.
\end{proof}

\begin{corollary}[Graph-universal leverage bound]
\label{cor:boundary-leverage-one}
Every response leverage in \cref{eq:boundary-leverage} satisfies
\begin{equation}
 \chi_v^2\leq Q_{vv}-\alpha=\frac{1-\alpha}{2}<1.
 \label{eq:boundary-leverage-one}
\end{equation}
Consequently
\begin{equation}
 \frac{\left|Q_{vS}(\widetilde x_S-x_S)\right|}{\sqrt{d_v}}
 \leq\delta
 \label{eq:normalized-boundary-error}
\end{equation}
under \cref{eq:energy-error}, without estimating a coordinate-specific
leverage.
\end{corollary}

\begin{proof}
The scalar Schur complement is
$K_v:=Q_{vv}-Q_{vS}Q_{SS}^{-1}Q_{Sv}$.  Its variational characterization and
$Q\succeq\alpha I$ give
\begin{equation*}
 K_v
 =\min_y\begin{bmatrix}y\\1\end{bmatrix}^{\!T}
 Q_{S\cup\{v\},S\cup\{v\}}
 \begin{bmatrix}y\\1\end{bmatrix}
 \geq\alpha.
\end{equation*}
The shared normalized matrix has $Q_{vv}=(1+\alpha)/2$, proving the first
claim.  Apply \cref{lem:boundary-interval} and $d_v\geq1$ for the displayed
normalized error bound.
\end{proof}

If the lower endpoint exceeds the gate, admission is safe.  If the upper
endpoint is below the gate, rejection is safe.  Only intervals intersecting
the transition band require refinement.  The lemma proves correctness of one
queried interval; it does not bound the work needed to find all ambiguous or
violating vertices.  That reporting problem is the main open response theorem.

The graph-universal leverage is intentionally oblivious to where a correction
was generated.  A frontier correction has much more structure: its dual
source is supported on the newly admitted batch.  The next lemma converts that
support into a trace bound on all boundary sensitivities.

\begin{lemma}[Source-aware leverage packing]
\label{lem:source-aware-leverage-packing}
Fix a face $S$, an exterior set $W\subseteq V\setminus S$, and a nonempty
source set $B\subseteq S$.  Let $E_B$ be the coordinate embedding from
$\mathbb R^B$ into $\mathbb R^S$ and put
\begin{equation}
 H_B:=E_B^TQ_{SS}^{-1}E_B.
 \label{eq:source-green-block}
\end{equation}
For $v\in W$, define the source-aware leverage
\begin{equation}
 (\chi_v^{B})^2
 :=Q_{vS}Q_{SS}^{-1}E_BH_B^{-1}
       E_B^TQ_{SS}^{-1}Q_{Sv}.
 \label{eq:source-aware-leverage}
\end{equation}
Then every correction of the form
\begin{equation}
 \Delta=Q_{SS}^{-1}E_Bh
 \label{eq:source-supported-correction}
\end{equation}
satisfies
\begin{equation}
 |Q_{vS}\Delta|
 \leq\chi_v^B\|\Delta\|_{Q_{SS}},
 \qquad v\in W,
 \label{eq:source-aware-interval}
\end{equation}
and the complete exterior leverage ledger obeys
\begin{equation}
 \sum_{v\in W}(\chi_v^B)^2\leq |B|.
 \label{eq:source-aware-leverage-sum}
\end{equation}
Consequently, for arbitrary positive normalized interval margins $m_v$, the
set
\begin{equation*}
 \mathcal A
 :=\left\{v\in W:
   \frac{\chi_v^B\|\Delta\|_{Q_{SS}}}{\sqrt{d_v}}
   \geq m_v
 \right\}
\end{equation*}
has the a priori packing bound
\begin{equation}
 \sum_{v\in\mathcal A}d_vm_v^2
 \leq |B|\,\|\Delta\|_{Q_{SS}}^2.
 \label{eq:source-aware-ambiguous-packing}
\end{equation}
For a common margin $m$, this gives
$\vol(\mathcal A)\leq |B|\|\Delta\|_{Q_{SS}}^2/m^2$.
Every nested exact correction in
\cref{prop:dual-local-orthogonal-corrections} has
\cref{eq:source-supported-correction} on $S=S_{j+1}$ with $B=B_j$.
\end{lemma}

\begin{proof}
The Green block $H_B$ is positive definite.  Define
\begin{equation*}
 P_B:=Q_{SS}^{-1/2}E_BH_B^{-1}E_B^TQ_{SS}^{-1/2}.
\end{equation*}
It is the orthogonal projector onto
$\operatorname{range}(Q_{SS}^{-1/2}E_B)$ and has trace $|B|$.
Also put $R:=Q_{WS}Q_{SS}^{-1/2}$.  Schur complementation of the principal
matrix on $S\cup W$, followed by $Q\preceq I$, gives
\begin{equation}
 RR^T=Q_{WS}Q_{SS}^{-1}Q_{SW}\preceq Q_{WW}\preceq I.
 \label{eq:exterior-response-contraction}
\end{equation}
Thus $R^TR\preceq I$ as well.  The definition
\cref{eq:source-aware-leverage} is the $v$th diagonal of $RP_BR^T$, so
\begin{equation*}
 \sum_{v\in W}(\chi_v^B)^2
 =\operatorname{tr}(P_BR^TR)
 \leq\operatorname{tr}(P_B)=|B|,
\end{equation*}
proving \cref{eq:source-aware-leverage-sum}.

For \cref{eq:source-supported-correction}, let
$w_v:=E_B^TQ_{SS}^{-1}Q_{Sv}$.  Then
\begin{equation*}
 Q_{vS}\Delta=w_v^Th,
 \qquad
 \|\Delta\|_{Q_{SS}}^2=h^TH_Bh.
\end{equation*}
Cauchy--Schwarz in the $H_B$ inner product gives
$|w_v^Th|^2\leq(w_v^TH_B^{-1}w_v)(h^TH_Bh)$, which is exactly
\cref{eq:source-aware-interval}.  Every $v\in\mathcal A$ satisfies
$d_vm_v^2\leq(\chi_v^B)^2\|\Delta\|_{Q_{SS}}^2$; summing and applying
\cref{eq:source-aware-leverage-sum} proves
\cref{eq:source-aware-ambiguous-packing}.  Finally,
\cref{eq:dual-frontier-support} says that
$Q_{S_{j+1}S_{j+1}}\Delta^j$ is supported on $B_j$, which gives the last
claim by solving the enlarged principal system.
\end{proof}

The source-aware ledger is an a priori statement: it bounds the volume whose
certified interval can be wide at a declared margin before the actual exterior
change is computed.  The following identity shows that its masses are not new
matrix objects.  They are exactly the diagonal loss of the terminal Schur
operator across the corresponding elimination.

\begin{theorem}[Source-aware leverage equals Schur-diagonal loss]
\label{thm:source-leverage-diagonal-loss}
Let $A\subsetneq T=A\mathbin{\dot\cup}B$ and
$W\subseteq V\setminus T$.  Write
\begin{equation*}
 L^A:=Q/A,
 \qquad
 L^T:=Q/T
\end{equation*}
for the terminal Schur complements, restricted to their surviving
coordinates.  On the enlarged face $T$, let $E_B$ embed $\mathbb R^B$ into
$\mathbb R^T$, put $H_B:=E_B^TQ_{TT}^{-1}E_B$, and define
\begin{equation}
 (\chi_v^{B\mid A})^2
 :=Q_{vT}Q_{TT}^{-1}E_BH_B^{-1}
      E_B^TQ_{TT}^{-1}Q_{Tv},
 \qquad v\in W.
 \label{eq:conditional-source-leverage}
\end{equation}
Then
\begin{equation}
 (\chi_v^{B\mid A})^2
 =L^A_{vB}(L^A_{BB})^{-1}L^A_{Bv}
 =L^A_{vv}-L^T_{vv}.
 \label{eq:source-leverage-diagonal-loss}
\end{equation}
Consequently every group $C\subseteq W$ has the exact additive mass
\begin{equation}
 Z_{B\mid A}(C)
 :=\sum_{v\in C}\frac{(\chi_v^{B\mid A})^2}{d_v}
 =\operatorname{tr}\!\left[
  D_C^{-1}\bigl(L^A_{CC}-L^T_{CC}\bigr)
 \right].
 \label{eq:group-diagonal-loss}
\end{equation}

More generally, let
$S_0\subsetneq\cdots\subsetneq S_J$ and
$B_j:=S_{j+1}\setminus S_j$.  If $v\notin S_k$, then its source-aware
sensitivity has the temporal telescope
\begin{equation}
 \sum_{j=0}^{k-1}(\chi_v^{B_j\mid S_j})^2
 =(Q/S_0)_{vv}-(Q/S_k)_{vv}
 \leq Q_{vv}-\alpha
 =\frac{1-\alpha}{2}.
 \label{eq:source-leverage-temporal-telescope}
\end{equation}
For every fixed $C\subseteq V\setminus S_k$, the normalized group masses
telescope as well:
\begin{equation}
 \sum_{j=0}^{k-1}Z_{B_j\mid S_j}(C)
 =\operatorname{tr}\!\left[
 D_C^{-1}\bigl((Q/S_0)_{CC}-(Q/S_k)_{CC}\bigr)
 \right].
 \label{eq:group-diagonal-temporal-telescope}
\end{equation}
Thus every still-exterior label has a graph-universal total squared
source-aware sensitivity budget, independent of the number of preceding
expansions.
\end{theorem}

\begin{proof}
Order $T$ as $(A,B)$ and put $K_B:=L^A_{BB}$.  The block inverse formula gives
\begin{equation*}
 H_B=K_B^{-1},
 \qquad
 Q_{vT}Q_{TT}^{-1}E_B=L^A_{vB}K_B^{-1}.
\end{equation*}
Substitution in \cref{eq:conditional-source-leverage} proves the first
equality in \cref{eq:source-leverage-diagonal-loss}.  Schur associativity
gives
\begin{equation*}
 L^T_{WW}
 =L^A_{WW}-L^A_{WB}K_B^{-1}L^A_{BW},
\end{equation*}
whose $v$th diagonal proves the second equality.  Summing these diagonal
identities with weights $1/d_v$ proves \cref{eq:group-diagonal-loss}.

Apply the same identity at every consecutive pair $(S_j,S_{j+1})$ and
telescope.  Since $Q\succeq\alpha I$, the variational characterization of a
Schur complement gives $Q/S_k\succeq\alpha I$.  Also
$(Q/S_0)_{vv}\leq Q_{vv}$.  The shared PageRank matrix has
$Q_{vv}=(1+\alpha)/2$, proving
\cref{eq:source-leverage-temporal-telescope}.  Summing the same termwise
diagonal telescope over $v\in C$ with weights $1/d_v$ proves
\cref{eq:group-diagonal-temporal-telescope}.
\end{proof}

The identity replaces source-conditioned response matrices by one monotone
scalar field: terminal Schur diagonals.  It does not make those diagonals
free, but a dynamic implementation may now obtain every group mass by a
before--after aggregate diagonal query.  In particular, no multiplicative
accuracy is needed when the true decrement is far below the current search
scale.  The next lemma first states the clean multiplicative hierarchy, and
the following corollary weakens its oracle to a one-sided additive guarantee.

\begin{lemma}[Rank-sensitive hierarchical locator]
\label{lem:source-aware-hierarchical-locator}
Under the hypotheses of
\cref{lem:source-aware-leverage-packing}, let $\mathcal H$ be a rooted
partition hierarchy whose leaves are the vertices of $W$, whose arity is at
most $b$, and whose depth is $L$.  For every node $C\subseteq W$, define its
degree-normalized source-aware mass by
\begin{equation}
 Z_B(C):=\sum_{v\in C}\frac{(\chi_v^B)^2}{d_v}.
 \label{eq:source-aware-group-mass}
\end{equation}
Suppose a group query returns a certified multiplicative upper estimate
\begin{equation}
 Z_B(C)\leq\widehat Z_B(C)\leq\gamma Z_B(C)
 \label{eq:source-aware-group-oracle}
\end{equation}
for some $\gamma\geq1$.  Fix a correction
$\Delta=Q_{SS}^{-1}E_Bh$ with energy radius
$E:=\|\Delta\|_{Q_{SS}}>0$ and a common normalized margin $m>0$.

Starting at the root, discard a queried node $C$ when
$\widehat Z_B(C)<m^2/E^2$ and otherwise query its children.  At every reached
leaf, evaluate the exact leaf predicate.  This search returns precisely the
potentially ambiguous set
\begin{equation*}
 \mathcal A_m
 :=\left\{v\in W:
     \frac{\chi_v^B E}{\sqrt{d_v}}\geq m
   \right\},
\end{equation*}
uses at most
\begin{equation}
 1+bL\left(1+\gamma|B|\frac{E^2}{m^2}\right)
 \label{eq:source-aware-group-query-count}
\end{equation}
group queries, and has output charge
\begin{equation}
 \vol(\mathcal A_m)\leq |B|\frac{E^2}{m^2}.
 \label{eq:source-aware-locator-output}
\end{equation}
Thus, if group queries and exact reached-leaf tests cost polylogarithmic work,
the search is sensitive to the packed ambiguity budget rather than to the
ambient boundary size, up to the hierarchy depth and approximation factor.
\end{lemma}

\begin{proof}
The masses in \cref{eq:source-aware-group-mass} are nonnegative and additive
over the children of every hierarchy node.  Since $d_v\geq1$,
\cref{eq:source-aware-leverage-sum} gives
\begin{equation}
 Z_B(W)\leq\sum_{v\in W}(\chi_v^B)^2\leq |B|.
 \label{eq:source-aware-normalized-mass-total}
\end{equation}
If $v\in\mathcal A_m$, then every ancestor $C$ of $v$ satisfies
$Z_B(C)\geq Z_B(\{v\})\geq m^2/E^2$, and therefore cannot be discarded by
the upper estimate in \cref{eq:source-aware-group-oracle}.  Conversely, the
exact predicate at a reached leaf removes every false candidate, proving
correctness.

At any fixed depth the retained nodes are disjoint.  A retained node satisfies
$\widehat Z_B(C)\geq m^2/E^2$, and the upper side of
\cref{eq:source-aware-group-oracle} implies
$Z_B(C)\geq m^2/(\gamma E^2)$.  Hence
\cref{eq:source-aware-normalized-mass-total} bounds the number of retained
nodes at that depth by $\gamma|B|E^2/m^2$.  Querying at most $b$ children per
retained node over $L$ levels, together with the root and harmless integer
rounding, proves \cref{eq:source-aware-group-query-count}.  The output-volume
claim is the common-margin case of
\cref{eq:source-aware-ambiguous-packing}.
\end{proof}

\begin{corollary}[Margin-adaptive additive group oracle]
\label{cor:additive-group-locator}
In \cref{lem:source-aware-hierarchical-locator}, put
$\tau:=m^2/E^2$ and replace the multiplicative oracle by certified estimates
satisfying
\begin{equation}
 Z_B(C)\leq\widehat Z_B(C)
 \leq\gamma Z_B(C)+\beta\tau,
 \qquad 0\leq\beta<1.
 \label{eq:additive-group-oracle}
\end{equation}
Use the same rule: discard $C$ when $\widehat Z_B(C)<\tau$, and otherwise
descend.  The search still returns exactly $\mathcal A_m$, and it uses at
most
\begin{equation}
 1+bL\left(
  1+\frac{\gamma}{1-\beta}|B|\frac{E^2}{m^2}
 \right)
 \label{eq:additive-group-query-count}
\end{equation}
group queries.  The oracle may therefore have additive error proportional to
the current pruning threshold; it need not resolve arbitrarily small group
masses multiplicatively.  By
\cref{thm:source-leverage-diagonal-loss}, the same statement holds when each
$Z_B(C)$ is supplied as the weighted terminal-Schur diagonal decrement in
\cref{eq:group-diagonal-loss}.
\end{corollary}

\begin{proof}
If $v\in\mathcal A_m$, every ancestor has
$Z_B(C)\geq Z_B(\{v\})\geq\tau$, so the certified upper estimate cannot
discard it.  Exact leaf tests again remove all false candidates.  Conversely,
every retained node satisfies
\begin{equation*}
 \tau\leq\widehat Z_B(C)
 \leq\gamma Z_B(C)+\beta\tau,
\end{equation*}
and hence $Z_B(C)\geq(1-\beta)\tau/\gamma$.  Retained nodes at one depth are
disjoint, while $Z_B(W)\leq|B|$.  The level count and the child-query charge
are therefore the same as in
\cref{lem:source-aware-hierarchical-locator}, with the factor
$\gamma/(1-\beta)$, proving
\cref{eq:additive-group-query-count}.
\end{proof}

The group mass has the exact trace form
\begin{equation*}
 Z_B(C)
 =\operatorname{tr}\!\left(
 P_BR_C^TD_C^{-1}R_C
 \right),
 \qquad
 R_C:=Q_{CS}Q_{SS}^{-1/2},
\end{equation*}
where $D_C:=\operatorname{diag}(d_v:v\in C)$.  Thus the remaining P0
primitive is no longer an unspecified boundary scan: it is a dynamic,
constant-factor group-trace oracle for these source-conditioned response
rows, together with exact evaluation only at the reached leaves.  The lemma
does not construct that oracle on a general nonequitable cyclic core.

The next reduction shows that the trace oracle does not require one response
query per hierarchy node.  A logarithmic batch of source-normalized random
responses estimates every fixed node simultaneously, including nodes chosen
adaptively by the subsequent branch-and-bound search.

\begin{lemma}[Group traces from logarithmically many response sketches]
\label{lem:group-trace-response-sketch}
Under the hypotheses of
\cref{lem:source-aware-hierarchical-locator}, let $\mathcal C$ be the set of
$M$ nodes of the fixed hierarchy.  Fix $\xi,\delta\in(0,1)$, draw
$r$ independent standard Gaussian vectors
$g_1,\ldots,g_r\in\mathbb R^B$, and define the exterior response sketches
\begin{equation}
 s^{(t)}
 :=D_W^{-1/2}Q_{WS}Q_{SS}^{-1}E_BH_B^{-1/2}g_t,
 \qquad 1\leq t\leq r.
 \label{eq:source-normalized-response-sketch}
\end{equation}
For every hierarchy node $C$, put
\begin{equation}
 \overline Z_B(C)
 :=\frac1r\sum_{t=1}^r\|s_C^{(t)}\|_2^2,
 \qquad
 \widehat Z_B(C):=\frac{\overline Z_B(C)}{1-\xi}.
 \label{eq:sketched-group-trace}
\end{equation}
There is a universal constant $c_{\rm sk}>0$ such that, if
\begin{equation}
 r\geq c_{\rm sk}\xi^{-2}\log(2M/\delta),
 \label{eq:group-trace-sketch-count}
\end{equation}
then, with probability at least $1-\delta$, simultaneously for every
$C\in\mathcal C$,
\begin{equation}
 Z_B(C)\leq\widehat Z_B(C)
 \leq\frac{1+\xi}{1-\xi}Z_B(C).
 \label{eq:sketched-group-oracle}
\end{equation}
In particular, $\xi=1/3$ supplies the oracle in
\cref{eq:source-aware-group-oracle} with $\gamma=2$ using
$r=O(\log(M/\delta))$ response sketches.
\end{lemma}

\begin{proof}
For a node $C$, define
\begin{equation*}
 T_C:=D_C^{-1/2}Q_{CS}Q_{SS}^{-1}E_BH_B^{-1/2}.
\end{equation*}
Cyclicity of trace and the definition of $P_B$ give
\begin{equation*}
 \|T_C\|_F^2
 =\operatorname{tr}\!\left(P_BR_C^TD_C^{-1}R_C\right)
 =Z_B(C).
\end{equation*}
Moreover, $s_C^{(t)}=T_Cg_t$, so
$\overline Z_B(C)=r^{-1}\sum_tg_t^T(T_C^TT_C)g_t$.  The standard Gaussian
quadratic-form moment-generating-function bound gives, for every positive
semidefinite matrix $A$ and $0<\xi<1$,
\begin{equation*}
 \Pr\!\left[
  \left|\frac1r\sum_{t=1}^rg_t^TAg_t-\operatorname{tr}(A)\right|
  >\xi\operatorname{tr}(A)
 \right]
 \leq2\exp(-c r\xi^2)
\end{equation*}
for a universal $c>0$.  Apply this with $A=T_C^TT_C$ and take a union bound
over all $M$ hierarchy nodes.  On the resulting event,
$(1-\xi)Z_B(C)\leq\overline Z_B(C)\leq
(1+\xi)Z_B(C)$ simultaneously.  Rescaling the estimate proves
\cref{eq:sketched-group-oracle}.  Nodes of zero mass have $T_C=0$ and obey
the same statement deterministically.
\end{proof}

Once the leaf values of the sketches are available, all quantities in
\cref{eq:sketched-group-trace} are additive and can be aggregated bottom-up.
The remaining challenge is therefore narrower still: maintain a polylogarithmic
batch of the signed response vectors in
\cref{eq:source-normalized-response-sketch}, and their squared subtree sums,
under terminal additions without materializing every exterior coordinate.
The lemma is a static randomized reduction and does not claim that dynamic
maintenance has local work.

Orthogonal frontier lifts make the sketches causal and mergeable across an
adaptive support trace.  The important choice is to accumulate the
nonnegative per-batch squared masses, rather than squaring one signed sketch
that has been reused to choose later batches.  Fresh randomness after each
batch avoids an adaptive-sketching gap.

\begin{theorem}[Causal append-only orthogonal response sketches]
\label{thm:causal-orthogonal-response-sketches}
Use the nested faces, Schur complements, and lifts of
\cref{thm:orthogonal-frontier-lifts}, and define the normalized lifted basis
\begin{equation}
 U_j:=\mathcal L_jK_j^{-1/2}\in\mathbb R^{V\times B_j}.
 \label{eq:normalized-frontier-basis}
\end{equation}
The bordered inverse identity also gives
\begin{equation}
 E_{B_j}^TQ_{S_{j+1}S_{j+1}}^{-1}E_{B_j}=K_j^{-1},
 \qquad
 U_j
 =Q_{S_{j+1}S_{j+1}}^{-1}E_{B_j}K_j^{1/2},
 \label{eq:normalized-frontier-source-identity}
\end{equation}
where the second identity is extended by zero outside $S_{j+1}$.  Thus its
exterior row norms are exactly the source-aware leverages of
\cref{eq:source-aware-leverage}.
Then
\begin{equation}
 U_i^TQU_j=
 \begin{cases}
 I_{B_j},&i=j,\\
 0,&i\neq j.
 \end{cases}
 \label{eq:orthonormal-frontier-bases}
\end{equation}
For an exterior vertex $v\notin S_{k+1}$, put
\begin{equation}
 \lambda_{jv}^2
 :=\frac{\|e_v^TQU_j\|_2^2}{d_v},
 \qquad
 \Lambda_{kv}^2:=\sum_{j=0}^k\lambda_{jv}^2.
 \label{eq:cumulative-frontier-leverage}
\end{equation}
If $Q\preceq I$ and $R_k:=\sum_{j=0}^k|B_j|$, then
\begin{equation}
 \sum_{v\notin S_{k+1}}\Lambda_{kv}^2\leq R_k.
 \label{eq:cumulative-frontier-leverage-packing}
\end{equation}
Moreover, every accumulated lifted error
$e=\sum_{j=0}^kU_ja_j$ satisfies
\begin{equation}
 \frac{|(Qe)_v|}{\sqrt{d_v}}
 \leq\Lambda_{kv}\|e\|_Q.
 \label{eq:cumulative-frontier-interval}
\end{equation}

Now allow the exact support batches and current boundary hierarchy to depend
on every earlier random choice.  After $B_j$, $K_j$, and the current hierarchy
are fixed, draw fresh independent standard Gaussian vectors
$g_{j,1},\ldots,g_{j,r_j}\in\mathbb R^{B_j}$ and define the per-batch group
increment
\begin{equation}
 \overline Z_j(C)
 :=\frac1{r_j}\sum_{t=1}^{r_j}
 \left\|D_C^{-1/2}Q_{CV}U_jg_{j,t}\right\|_2^2.
 \label{eq:append-only-group-mass-increment}
\end{equation}
Let $M_j$ bound the number of live hierarchy nodes and leaves at event $j$,
fix $\xi,\delta\in(0,1)$, and set
\begin{equation}
 \delta_j:=\frac{6\delta}{\pi^2(j+1)^2}.
 \label{eq:adaptive-sketch-failure-budget}
\end{equation}
For
\begin{equation}
 r_j\geq c_{\rm sk}\xi^{-2}\log(2M_j/\delta_j),
 \label{eq:adaptive-sketch-count}
\end{equation}
with probability at least $1-\delta$, simultaneously at every event, every
stored leaf and every current group has cumulative estimate
\begin{equation}
 \sum_{i\leq j}Z_i(C)
 \leq
 \frac1{1-\xi}\sum_{i\leq j}\overline Z_i(C)
 \leq
 \frac{1+\xi}{1-\xi}\sum_{i\leq j}Z_i(C),
 \label{eq:adaptive-cumulative-group-oracle}
\end{equation}
where
$Z_i(C):=\|D_C^{-1/2}Q_{CV}U_i\|_F^2$ and contributions are retained only
while their leaf remains exterior.  If a vertex $v$ first enters the boundary
after batch $B_j$, then $\lambda_{iv}=0$ for every $i<j$.  Its accumulator
therefore starts at zero and no earlier response sketch is replayed.  Put
$N:=1+\operatorname{cvol}(S_J)$.  Since every batch is nonempty,
$J\leq |S_J|$; if $M_j\leq N^{O(1)}$, the total number of sampled response
right-hand sides is $O_\xi(J\log(NJ/\delta))$.
\end{theorem}

\begin{proof}
The bottom-right block of the bordered inverse is $K_j^{-1}$, proving the
first identity in \cref{eq:normalized-frontier-source-identity}.  Also,
$Q_{S_{j+1}S_{j+1}}\mathcal L_jy=E_{B_j}K_jy$; substitute
$y=K_j^{-1/2}a$ and solve the principal system to obtain the second identity.
The isometry and cross-space orthogonality in
\cref{eq:frontier-lift-isometry,eq:frontier-lift-orthogonality} give
\cref{eq:orthonormal-frontier-bases}.  Let
$U_{\leq k}:=[U_0\ \cdots\ U_k]$.  Then
$Q^{1/2}U_{\leq k}$ has $R_k$ orthonormal columns.  For
$W:=V\setminus S_{k+1}$,
\begin{align*}
 \sum_{v\in W}\Lambda_{kv}^2
 &=\|D_W^{-1/2}Q_{WV}U_{\leq k}\|_F^2\\
 &\leq\|Q_{WV}U_{\leq k}\|_F^2
 \leq R_k,
\end{align*}
because $d_v\geq1$ and
$\|E_W^TQ^{1/2}\|_2\leq\|Q^{1/2}\|_2\leq1$.  Writing
$a:=(a_0,\ldots,a_k)$, orthonormality gives
$\|e\|_Q=\|a\|_2$; rowwise Cauchy--Schwarz proves
\cref{eq:cumulative-frontier-interval}.

At event $j$, condition on the complete history through the choice of
$U_j$ and the current hierarchy.  They are fixed before the fresh Gaussian
vectors are drawn.  Apply the quadratic-form concentration argument from
\cref{lem:group-trace-response-sketch} with failure probability $\delta_j$
to all current nodes and leaves.  Thus, conditionally on every possible
history, the event-$j$ estimates lie between
$(1-\xi)Z_j(C)$ and $(1+\xi)Z_j(C)$ except with probability $\delta_j$.
Since $\sum_{j\geq0}\delta_j=\delta$, a conditional union bound proves the
simultaneous claim despite adaptive future batches.  Summing nonnegative
per-batch estimates preserves both multiplicative inequalities; group
estimates are sums of their retained leaf estimates.

Finally, suppose $v$ first becomes adjacent to the active face after $B_j$.
For $i<j$, the support of $U_i$ lies in
$S_{i+1}\subseteq S_j$, while $v$ has no edge to $S_j$.  Hence
$e_v^TQU_i=0$, proving causal zero initialization.
Finally, $J\leq|S_J|$ follows from disjoint nonempty batches.  Under the
polynomial hierarchy-size assumption,
\cref{eq:adaptive-sketch-count,eq:adaptive-sketch-failure-budget} gives
$r_j=O_\xi(\log(NJ/\delta))$ uniformly; summing over the events proves the
sample count.
\end{proof}

\begin{corollary}[Cumulative sketches collapse to diagonal loss]
\label{cor:cumulative-sketch-diagonal-loss}
In \cref{thm:causal-orthogonal-response-sketches}, fix $k$ and a hierarchy
node $C\subseteq V\setminus S_{k+1}$.  The exact cumulative mass being
estimated has the deterministic representation
\begin{equation}
 \sum_{j=0}^k Z_j(C)
 =\operatorname{tr}\!\left[
 D_C^{-1}\bigl((Q/S_0)_{CC}-(Q/S_{k+1})_{CC}\bigr)
 \right].
 \label{eq:cumulative-sketch-diagonal-loss}
\end{equation}
Therefore a dynamic oracle that directly supplies the one-sided additive
estimates in \cref{eq:additive-group-oracle} for these diagonal losses can
replace all per-batch Gaussian response sketches.  Combined with
\cref{cor:additive-group-locator}, it gives the same no-false-negative
branch-and-bound search with output-sensitive query count.  No source vector,
old sketch, or past batch must be replayed.
\end{corollary}

\begin{proof}
The normalized lifted basis $U_j$ is the source-normalized response basis for
$B_j$ by \cref{eq:normalized-frontier-source-identity}.  Hence
\begin{equation*}
 Z_j(C)
 =\sum_{v\in C}\frac{(\chi_v^{B_j\mid S_j})^2}{d_v}.
\end{equation*}
Apply \cref{eq:group-diagonal-temporal-telescope} with $k+1$ in place of
$k$.  The oracle and search claim is then exactly
\cref{cor:additive-group-locator}.
\end{proof}

\begin{theorem}[Static all-boundary checkpoint at accelerated cost]
\label{thm:static-diagonal-loss-checkpoint}
For an exposed face $S$, let
$W_S:=\{v\notin S:Q_{vS}\neq0\}$ and define
\begin{equation}
 h_v(S):=\frac{Q_{vS}Q_{SS}^{-1}Q_{Sv}}{d_v}
 =\frac{Q_{vv}-(Q/S)_{vv}}{d_v},
 \qquad v\in W_S.
 \label{eq:static-checkpoint-loss}
\end{equation}
Fix $\zeta,\xi,\delta\in(0,1/2)$.  A degree
\begin{equation*}
 m=O\!\left(\alpha^{-1/2}\log\frac1{\alpha\zeta}\right)
\end{equation*}
Chebyshev polynomial $p_m$ and
\begin{equation*}
 r=O\!\left(\xi^{-2}\log\frac{2\max\{1,|W_S|\}}{\delta}\right)
\end{equation*}
independent standard Gaussian vectors $g_t\in\mathbb R^S$ give the estimates
\begin{equation}
 \overline h_v:=\frac1r\sum_{t=1}^r
  \frac{\bigl(Q_{vS}p_m(Q_{SS})g_t\bigr)^2}{d_v},
 \qquad
 \widehat h_v:=\frac{\overline h_v}
 {(1-\xi)(1-\zeta)^2}.
 \label{eq:static-checkpoint-estimator}
\end{equation}
With probability at least $1-\delta$, simultaneously for every $v\in W_S$,
\begin{equation}
 h_v(S)\leq\widehat h_v
 \leq
 \frac{(1+\xi)(1+\zeta)^2}
 {(1-\xi)(1-\zeta)^2}h_v(S).
 \label{eq:static-checkpoint-guarantee}
\end{equation}
The coordinate estimates, and therefore every group sum in a fixed boundary
hierarchy, are constructed in
\begin{equation}
 \widetilde O_{\zeta,\xi,\delta}\!\left(
  \frac{\operatorname{cvol}(S)}{\sqrt\alpha}
 \right)
 \label{eq:static-checkpoint-work}
\end{equation}
work from the rows of $S$ and the degrees of their incident boundary labels.
\end{theorem}

\begin{proof}
The spectrum of the principal matrix $Q_{SS}$ lies in $[\alpha,1]$.
Chebyshev approximation on this interval supplies $p_m$ with
\begin{equation*}
 |p_m(x)-x^{-1/2}|\leq\zeta x^{-1/2},
 \qquad x\in[\alpha,1],
\end{equation*}
at the stated degree.  Indeed, after mapping the interval to $[-1,1]$, the
nearest singularity has limiting Bernstein-ellipse parameter
$(1+\sqrt\alpha)/(1-\sqrt\alpha)$.  Therefore
\begin{equation*}
 (1-\zeta)^2Q_{SS}^{-1}
 \preceq p_m(Q_{SS})^2
 \preceq(1+\zeta)^2Q_{SS}^{-1}.
\end{equation*}
For each $v$, the summand before squaring in
\cref{eq:static-checkpoint-estimator} is centered Gaussian with variance
\begin{equation*}
 \mu_v:=d_v^{-1}Q_{vS}p_m(Q_{SS})^2Q_{Sv},
\end{equation*}
which lies between $(1-\zeta)^2h_v(S)$ and
$(1+\zeta)^2h_v(S)$.  Moreover
$\overline h_v/\mu_v\sim\chi_r^2/r$ when $\mu_v>0$, while a zero row is
exactly zero.  Chi-square concentration and a union bound prove
\cref{eq:static-checkpoint-guarantee}.  Summing the coordinate inequalities
gives all group inequalities.

A Chebyshev recurrence applies $p_m(Q_{SS})$ to one vector using $m$ row
scans of $S$.  One additional cut-incidence scan forms all
$Q_{W_S,S}p_m(Q_{SS})g_t$ simultaneously.  Repeating for the $r$ probes proves
\cref{eq:static-checkpoint-work}.
\end{proof}

\begin{corollary}[Geometric checkpoint amortization]
\label{cor:static-diagonal-geometric-amortization}
If checkpoint faces satisfy
$\operatorname{cvol}(S_{i+1})\geq
(1+\gamma)\operatorname{cvol}(S_i)$ for fixed $\gamma>0$, then independent
instances of \cref{thm:static-diagonal-loss-checkpoint}, with summable failure
budgets, construct every full boundary snapshot in total work
\begin{equation}
 \widetilde O\!\left(
  \frac{\operatorname{cvol}(S_K)}{\sqrt\alpha}
 \right)
 \label{eq:static-diagonal-geometric-work}
\end{equation}
and fail with at most the sum of their budgets.
\end{corollary}

\begin{proof}
Sum \cref{eq:static-checkpoint-work} and use
$\sum_i\operatorname{cvol}(S_i)\leq
(1+1/\gamma)\operatorname{cvol}(S_K)$.  A union bound gives the failure
claim.
\end{proof}

This closes the Monte Carlo work bound for full-boundary state at anchor
checkpoints; deterministic or Las Vegas certification is separate.  It does
not close the epoch-internal problem: subtracting two noisy cumulative
snapshots cannot certify a much smaller intervening loss.  The remaining
oracle must therefore maintain the incremental post-checkpoint loss with
scale-matched additive error, or obtain the equivalent scalar harmonic
measurements only for the nodes reached by the packed search.
